\begin{tikzpicture}[x=1cm,y=-.87cm,every node/.style={font=\fontsize{12.2}{15}\selectfont,text=Ink},arr/.style={-{Stealth[length=2mm]},line width=.8pt,draw=Classical}]\path[use as bounding box] (0,0) rectangle (7.45,6.2);
\draw[draw=Classical,fill=Classical!8,rounded corners=2pt,line width=.7pt] (0.6,0.7) rectangle (6.8,2.2);
\node[anchor=center,text=Ink,font=\fontsize{12.2}{15}\selectfont,text width=6cm,align=center] at (3.7,1.45) {Sezione 2D: area dei pori};
\draw[arr,draw=Quantum] (3.7,2.4)--(3.7,3.2);
\draw[draw=Quantum,fill=Quantum!8,rounded corners=2pt,line width=.7pt] (0.6,3.4) rectangle (6.8,4.9);
\node[anchor=center,text=Ink,font=\fontsize{12.2}{15}\selectfont,text width=6cm,align=center] at (3.7,4.15) {Carota: permeabilità al fluido};
\node[anchor=center,text=Ink,font=\fontsize{11.5}{15}\selectfont,text width=6.5cm,align=center] at (3.7,5.7) {Scale e proprietà fisiche differenti};

\end{tikzpicture}